%% ============================================================================
%%  sobre2-comun.tex -- Datos y utilidades comunes del Sobre N.° 2
%%
%%  Lo carga la hoja resumen del sobre (LAFROX-Sobre2-Indice.tex) después de
%%  \documentclass. Los subdocumentos se compilan desde sus propias ramas;
%%  compilar_sobre2.py les fija el folio inicial sin modificar sus fuentes.
%% ============================================================================

\datoslafrox{
  caso       = {Caso --- Logística},
  fecha      = {25 de noviembre de 2026},
  documento  = {Propuesta Técnica},
  sobre      = {Sobre N.\textdegree{} 2 \textemdash{} Oferta Técnica},
}

% El número de página se rotula como folio, igual que en el Sobre N.° 1.
\renewcommand{\LFXfolio}{Folio}

% Art. 40.2: la portada también lleva media firma, además de la firma completa.
\renewcommand{\LFXpiefolioen}[1]{\LFXpiederecho}

% ---- Foliación correlativa del sobre ---------------------------------------
%  folios-sobre2.tex lo escribe compilar_sobre2.py: trae las filas de la hoja
%  resumen con el folio de inicio y de término de cada documento.
\newcommand{\filasindicesobre}{}
\InputIfFileExists{folios-sobre2.tex}{}{}
\providecommand{\LFXfolioinicial}{1}

% ---- Tabla cuyas filas vienen en una macro ---------------------------------
%  tablalafrox lee su cuerpo como argumento y le agrega los filetes de fila;
%  para que vea filas reales, la macro de filas se expande una vez antes.
%  \tablaconcuerpo{caption}{label}{columnas}{cabecera}{\macrodefilas}
\makeatletter
\newcommand{\tablaconcuerpo}[5]{%
  \expandafter\LFX@tablaconcuerpo\expandafter{#5}{#1}{#2}{#3}{#4}}
\newcommand{\LFX@tablaconcuerpo}[5]{%
  \begin{tablalafrox}{#2}{#3}{#4}{#5}#1\end{tablalafrox}}
\makeatother
